\addcontentsline{toc}{chapter}{СПИСОК ИСПОЛЬЗОВАННЫХ ИСТОЧНИКОВ}
\begin{thebibliography}{}
    \bibitem{volodina} Володина Е.В., Студентова Е.А. Практическое применение алгоритма решения задачи коммивояжёра // Инженерный вестник Дона. - 2015. - №2.
    \bibitem{voronin} Воронин Е.В., Бутузова Е.А. ПРИМЕНЕНИЕ ЗАДАЧИ КОММИВОЯЖЁРА // Символ науки. - Уфа: Символ науки, 2020. - С. 62-63.
    \bibitem{sovet} Математический энциклопедический словарь.. - М.: Советская энциклопедия, 1988. - 848 с.
    \bibitem{schrijver} Alexander Schrijver. On the history of combinatorial optimization (till 1960) // 2005
    \bibitem{mudrov} Мудров В.И. Задача о Коммивояжёре. - М.: Знание, 1969. - 64 с.
    \bibitem{belousov} Белоусов А.И., Ткачев С.Б. Дискретная математика: учебник для вузов. - 7-е изд. - М.: Издательство МГТУ им. Н. Э. Баумана, 2021. - 703 с.
    \bibitem{asa} Асанов М.О., Ткачев С.Б. Дискретная математика: графы, матроиды, алгоритмы : учебное пособие. - 3-е изд. - СПб.: Лань, 2020. - 364 с
    \bibitem{BRE} Задача коммивояжёра // Большая Российская Энциклопедия URL: https://bigenc.ru/c/zadachakommivoiazhera-tsp-61952a (дата обращения: 09.12.2024).
    \bibitem{tsp} Michael Junger, Gerhard Reinelt, Giovanni Rinaldi. The Traveling Salesman Problem. - Institut fur Informatik UNIVERSITAT ZU KOLN, 1994. - 112 с.
    \bibitem{her} Marco Dorigo, Thomas Stützle. Ant Colony Optimization: Overview and Recent Advances // Handbook of Metaheuristics. - Switzerland: Springer International Publishin, 2019. - С. 311-351.
    \bibitem{perms} Robert Sedgevick. Permutations generation methods // Computing Survey. - 1977
    \bibitem{perms-methods} Robert Sedgevick. Permutations generation methods // Princeton University. - 2002
\end{thebibliography}